\chapter{3-8译码器}

\section{第一步：解决什么问题}

译码器是编码器的逆过程。

编码器：8个输入 $\to$ 3位输出（压缩）

译码器：3位输入 $\to$ 8个输出（展开）

想象一个场景：处理器用3根地址线选择8个外设中的一个。当处理器输出地址"101"时，第5号外设被选中（对应输出线变高），其余外设保持非选中状态。

\begin{keypoint}[译码器的本质]
译码器 = 信号展开器：$n$ 位输入 $\to$ $2^n$ 个输出

3-8译码器：3位输入 → 8个输出，恰好一个输出被激活
\end{keypoint}

\section{第二步：输入输出关系}

\begin{itemize}
  \item \textbf{输入}：$A_2, A_1, A_0$（3位地址输入）
  \item \textbf{输出}：$Y_0, Y_1, \dots, Y_7$（8个输出，通常高电平有效）
  \item \textbf{功能}：当输入 $(A_2 A_1 A_0) = k$ 时，$Y_k = 1$，所有其他输出为0
\end{itemize}

\section{第三步：真值表}

\begin{center}
\begin{tabular}{|c|c|c|c|c|c|c|c|c|c|}
\hline
$A_2$ & $A_1$ & $A_0$ & $Y_7$ & $Y_6$ & $Y_5$ & $Y_4$ & $Y_3$ & $Y_2$ & $Y_1$ & $Y_0$ \\
\hline
0 & 0 & 0 & 0 & 0 & 0 & 0 & 0 & 0 & 0 & 1 \\
0 & 0 & 1 & 0 & 0 & 0 & 0 & 0 & 0 & 1 & 0 \\
0 & 1 & 0 & 0 & 0 & 0 & 0 & 0 & 1 & 0 & 0 \\
0 & 1 & 1 & 0 & 0 & 0 & 0 & 1 & 0 & 0 & 0 \\
1 & 0 & 0 & 0 & 0 & 0 & 1 & 0 & 0 & 0 & 0 \\
1 & 0 & 1 & 0 & 0 & 1 & 0 & 0 & 0 & 0 & 0 \\
1 & 1 & 0 & 0 & 1 & 0 & 0 & 0 & 0 & 0 & 0 \\
1 & 1 & 1 & 1 & 0 & 0 & 0 & 0 & 0 & 0 & 0 \\
\hline
\end{tabular}
\end{center}

\section{第四步：逻辑推导}

译码器的每个输出对应输入的一个\textbf{最小项}：

\begin{align}
  Y_0 &= \overline{A_2} \cdot \overline{A_1} \cdot \overline{A_0} = m_0 \\
  Y_1 &= \overline{A_2} \cdot \overline{A_1} \cdot A_0 = m_1 \\
  Y_2 &= \overline{A_2} \cdot A_1 \cdot \overline{A_0} = m_2 \\
  Y_3 &= \overline{A_2} \cdot A_1 \cdot A_0 = m_3 \\
  Y_4 &= A_2 \cdot \overline{A_1} \cdot \overline{A_0} = m_4 \\
  Y_5 &= A_2 \cdot \overline{A_1} \cdot A_0 = m_5 \\
  Y_6 &= A_2 \cdot A_1 \cdot \overline{A_0} = m_6 \\
  Y_7 &= A_2 \cdot A_1 \cdot A_0 = m_7
\end{align}

\textbf{关键观察}：每个输出是一个3输入AND门！对于给定输入组合，恰好一个AND门的所有输入都为1（原变量或反变量），所以该AND门输出1。

\section{第五步：结构化设计}

译码器内部结构 = 8个3输入AND门 + 3个反相器：

\begin{center}
\begin{tikzpicture}[scale=0.9]
  % Input signals
  \node (A2) at (0,1) {$A_2$};
  \node (A1) at (0,0) {$A_1$};
  \node (A0) at (0,-1) {$A_0$};

  % Inverters
  \node[american not port] (inv2) at (1.5,1) {};
  \node[american not port] (inv1) at (1.5,0) {};
  \node[american not port] (inv0) at (1.5,-1) {};

  \draw (A2) -- (inv2.in);
  \draw (A1) -- (inv1.in);
  \draw (A0) -- (inv0.in);

  % AND gates
  \foreach \y in {0,...,7} {
    \node[american and port] (and\y) at (5, 1.5 - 0.5*\y) {};
    \node (Y\y) at (6.5, 1.5 - 0.5*\y) {$Y_\y$};
    \draw (and\y.out) -- (Y\y);
  }
\end{tikzpicture}
\end{center}

\textbf{设计规律}：每个AND门输入端根据该输出的二进制编号选择原变量或反变量。
\begin{itemize}
  \item $Y_5$（101）：$A_2$原变量 $\cdot$ $\overline{A_1}$反变量 $\cdot$ $A_0$原变量
  \item $Y_3$（011）：$\overline{A_2}$反变量 $\cdot$ $A_1$原变量 $\cdot$ $A_0$原变量
\end{itemize}

\textbf{一般规则}：对于 $Y_k$，如果 $k$ 的第 $i$ 位是0，则取 $\overline{A_i}$；如果第 $i$ 位是1，则取 $A_i$。

\section{第六步：为什么这样设计}

\begin{enumerate}
  \item \textbf{为什么每个输出恰好是一个AND门？}

  AND门的特性：\textbf{所有}输入必须为1，输出才为1。译码器的需求恰好是：对于输入组合$(A_2 A_1 A_0)$恰好等于某个特定值时，对应的输出线才为1。AND门天然实现了这个"完全匹配"检测。

  \item \textbf{为什么需要反相器？}

  因为输入信号可能是0。例如：当 $(A_2 A_1 A_0) = 000$ 时，要检测"$A_2=0$且$A_1=0$且$A_0=0$"，必须使用 $\overline{A_2} \cdot \overline{A_1} \cdot \overline{A_0}$ —— 需要三个反相器。

  \item \textbf{译码器与编码器的对偶关系？}

  编码器：输入线号 $\to$ 二进制编码（多对1的OR）
  译码器：二进制编码 $\to$ 输出线号（1对多的AND展开）

  这正是"编码"与"译码"的本质区别。
\end{enumerate}

\section{第七步：考试常见变式}

\begin{enumerate}
  \item \textbf{变式1：低电平有效译码器（如74LS138）}

  输入仍是高电平有效，但输出是低电平有效（选中时输出0，未选中时输出1）。此时将AND门换成NAND门即可，有时输出端还有小圆圈标记。

  74LS138还包含三个使能端：$G_1$（高有效），$\overline{G_{2A}}$和$\overline{G_{2B}}$（低有效）。\textbf{这几乎是必考题！}

  \item \textbf{变式2：用译码器实现任意逻辑函数}

  \textbf{核心考点}：译码器的输出是所有最小项。\textbf{任何组合逻辑函数}都可以表示为若干最小项的OR。

  例如：实现 $F(A,B,C) = \sum m(1,3,5,7)$（即 $F = C$）

  解法：将 $Y_1, Y_3, Y_5, Y_7$ 通过一个OR门连接即可。

  \item \textbf{变式3：用3-8译码器构建4-16译码器}

  使用两片3-8译码器 + 片选信号。高位地址作为片选，选择使用哪一个3-8译码器。

  \item \textbf{变式4：用译码器做地址译码}

  在存储器或外设选择中的应用。n位地址输入，$2^n$个片选信号输出。

  \item \textbf{变式5：显示译码器（BCD-7段译码器）}

  将4位BCD码译为7段数码管的驱动信号。每个输出段也是一个最小项组合的逻辑函数。
\end{enumerate}

\section{第八步：VHDL实现}

\subsection{普通3-8译码器（高电平有效）}

\begin{lstlisting}[caption={3-8译码器（case实现）}]
library ieee;
use ieee.std_logic_1164.all;

entity decoder_3to8 is
  port (
    A : in  std_logic_vector(2 downto 0);
    Y : out std_logic_vector(7 downto 0)
  );
end decoder_3to8;

architecture behavioral of decoder_3to8 is
begin
  process(A)
  begin
    case A is
      when "000" => Y <= "00000001";
      when "001" => Y <= "00000010";
      when "010" => Y <= "00000100";
      when "011" => Y <= "00001000";
      when "100" => Y <= "00010000";
      when "101" => Y <= "00100000";
      when "110" => Y <= "01000000";
      when "111" => Y <= "10000000";
      when others => Y <= "00000000";
    end case;
  end process;
end behavioral;
\end{lstlisting}

\subsection{74LS138风格（低电平有效 + 使能端）}

\begin{lstlisting}[caption={74LS138风格译码器}]
library ieee;
use ieee.std_logic_1164.all;

entity decoder_74ls138 is
  port (
    A  : in  std_logic_vector(2 downto 0);
    G1 : in  std_logic;
    G2A_n : in  std_logic;
    G2B_n : in  std_logic;
    Y_n : out std_logic_vector(7 downto 0)  -- 低电平有效
  );
end decoder_74ls138;

architecture behavioral of decoder_74ls138 is
  signal enable : std_logic;
begin
  enable <= G1 and (not G2A_n) and (not G2B_n);

  process(A, enable)
  begin
    if enable = '0' then
      Y_n <= "11111111";  -- 禁止：所有输出无效
    else
      case A is
        when "000" => Y_n <= "11111110";
        when "001" => Y_n <= "11111101";
        when "010" => Y_n <= "11111011";
        when "011" => Y_n <= "11110111";
        when "100" => Y_n <= "11101111";
        when "101" => Y_n <= "11011111";
        when "110" => Y_n <= "10111111";
        when "111" => Y_n <= "01111111";
        when others => Y_n <= "11111111";
      end case;
    end if;
  end process;
end behavioral;
\end{lstlisting}

\section{第九步：Testbench验证}

\begin{lstlisting}[caption={3-8译码器Testbench}]
library ieee;
use ieee.std_logic_1164.all;

entity decoder_3to8_tb is
end decoder_3to8_tb;

architecture test of decoder_3to8_tb is
  signal A : std_logic_vector(2 downto 0);
  signal Y : std_logic_vector(7 downto 0);

  component decoder_3to8 is
    port (
      A : in  std_logic_vector(2 downto 0);
      Y : out std_logic_vector(7 downto 0)
    );
  end component;
begin
  uut: decoder_3to8 port map (A, Y);

  stimulus: process
  begin
    for i in 0 to 7 loop
      A <= std_logic_vector(to_unsigned(i, 3));
      wait for 10 ns;
    end loop;
    wait;
  end process;
end test;
\end{lstlisting}

\section{第十步：如何快速解题}

\begin{examtip}[译码器快速解题法]
\begin{enumerate}
  \item \textbf{写输出表达式}：$Y_k = \prod A_i^{b_i}$，其中 $b_i$ 是 $k$ 的二进制第 $i$ 位（0取反，1取原）
  \item \textbf{用译码器实现逻辑函数}：将函数表示为最小项之和，将对应输出端用OR门连接
  \item \textbf{级联译码器}：高位地址 $\to$ 片选信号，低位地址 $\to$ 各片译码器的输入
  \item \textbf{有效电平}：高有效用AND，低有效用NAND
\end{enumerate}
\end{examtip}

\textbf{秒杀口诀}：每个输出 = 一个最小项 = 一个AND门（输入原/反变量取决于编号）。

\textbf{最常考：用3-8译码器 + OR门实现任意3变量逻辑函数。}

例子：实现 $F(A,B,C) = \sum m(0,2,4,6)$
解：$F = Y_0 + Y_2 + Y_4 + Y_6$（将四个输出通过OR门连接）

\section{变式详解}
\chapter{3-8译码器——5大变式详解}

\section{变式1：74LS138风格译码器（低有效输出+三使能端）——必考！}

\begin{detailbox}[完整教学]

\textbf{【电路】}
74LS138是考试中最常出现的译码器芯片。特点：
\begin{itemize}
  \item 3位地址输入A2,A1,A0
  \item 8个输出Y0-Y7，\textbf{全部低电平有效}（选中=0，未选中=1）
  \item 3个使能端：G1(高有效), G2A(低有效), G2B(低有效)
  \item 只有三个使能端同时有效（G1=1, G2A=0, G2B=0）时，译码器才工作
\end{itemize}

\textbf{【使能逻辑】}$enable = G1 \cdot \overline{G2A} \cdot \overline{G2B}$

\textbf{【VHDL】}
\begin{lstlisting}
entity decoder_74ls138 is
  port (
    A    : in  std_logic_vector(2 downto 0);
    G1   : in  std_logic;
    G2A_n: in  std_logic;  -- 低有效
    G2B_n: in  std_logic;  -- 低有效
    Y_n  : out std_logic_vector(7 downto 0)  -- 低有效
  );
end decoder_74ls138;

architecture behavioral of decoder_74ls138 is
  signal enable : std_logic;
begin
  enable <= G1 and (not G2A_n) and (not G2B_n);

  process(A, enable)
  begin
    if enable = '0' then
      Y_n <= "11111111";  -- 全高=全无效
    else
      case A is
        when "000" => Y_n <= "11111110";  -- Y0=0
        when "001" => Y_n <= "11111101";  -- Y1=0
        when "010" => Y_n <= "11111011";
        when "011" => Y_n <= "11110111";
        when "100" => Y_n <= "11101111";
        when "101" => Y_n <= "11011111";
        when "110" => Y_n <= "10111111";
        when "111" => Y_n <= "01111111";
        when others => Y_n <= "11111111";
      end case;
    end if;
  end process;
end behavioral;
\end{lstlisting}

\textbf{【为什么用3个使能端】}
多个使能端方便级联。G2A和G2B是低有效，G1是高有效——这种混合设计使得可以用不同信号控制，增加灵活性。

\textbf{【常考陷阱】}
\begin{itemize}
  \item 看清输出是低有效还是高有效——74LS138输出0表示选中
  \item 三个使能端必须\textbf{同时}满足条件译码器才工作
  \item 使能无效时所有输出为1（全高=全无效）
\end{itemize}

\end{detailbox}


\section{变式2：用译码器实现任意逻辑函数（必考！）}

\begin{detailbox}[完整教学]

\textbf{【核心原理】}
译码器的每个输出=一个最小项。\textbf{任何组合逻辑函数}=若干最小项之和（OR）。

\textbf{【通用公式】}$F(A,B,C) = \sum m(\dots)= Y_{i1}+Y_{i2}+Y_{i3}+\dots$

\textbf{【例题】}用3-8译码器+一个OR门实现$F(A,B,C)=\sum m(1,3,5,7)$。

\textbf{【解法】}
\begin{enumerate}
  \item 译码器产生8个最小项：Y0=m0, Y1=m1, ..., Y7=m7
  \item $F = m1+m3+m5+m7 = Y1+Y3+Y5+Y7$
  \item 用1个4输入OR门连接Y1,Y3,Y5,Y7
\end{enumerate}

\textbf{【电路结构】}
\begin{center}
\begin{tikzpicture}
\node[draw,minimum width=2.5cm,minimum height=1.5cm] (dec) at (0,0) {3-8译码器};
\node[american or port] (orgate) at (4,0) {};
\draw[->] (dec.east|-orgate.in 1) -- (orgate.in 1);
\draw[->] (dec.east|-orgate.in 2) -- (orgate.in 2);
\draw[->] (orgate.out) -- ++(1,0) node[right]{$F$};
\node at (0,-2) {A2A1A0 = 变量A,B,C};
\node at (4,-2) {OR门输入=Y1,Y3,Y5,Y7};
\end{tikzpicture}
\end{center}

\textbf{【VHDL】}
\begin{lstlisting}
-- 例化译码器
dec: decoder_3to8 port map(A=>ABC, Y=>Y);
-- OR门
F <= Y(1) or Y(3) or Y(5) or Y(7);
\end{lstlisting}

\textbf{【推广：实现任意3变量函数】}
\begin{center}
\begin{tabular}{|l|l|}
\hline
\textbf{函数} & \textbf{连接} \\
\hline
$F=m0+m2+m4+m6$ & Y0+Y2+Y4+Y6 \\
$F=m0+m1+m2+m3$ & Y0+Y1+Y2+Y3 \\
$F=A$ (即m4+m5+m6+m7) & Y4+Y5+Y6+Y7 \\
$F=\overline{A}$ & Y0+Y1+Y2+Y3 \\
\hline
\end{tabular}
\end{center}

\textbf{【如果函数用NAND实现】}
用NAND门替代OR门（德摩根定理）。此时输出是反相的，但功能等价。

\begin{examfocus}
\textbf{必考！用3-8译码器+门电路实现任意3变量逻辑函数。}
步骤：(1)将函数写成最小项之和 (2)对应Y输出用OR门连接
\end{examfocus}

\end{detailbox}


\section{变式3：用3-8译码器构建4-16译码器（片选级联）}

\begin{detailbox}[完整教学]

\textbf{【问题】}只有3-8译码器，如何得到4位地址→16个输出的4-16译码器？

\textbf{【方案】}两片3-8译码器+最高位地址A3做片选。

\textbf{【级联逻辑】}
\begin{itemize}
  \item \textbf{片0（低位）}：处理A3=0时，地址0000-0111（译码器输出Y0-Y7对应系统Y0-Y7）
  \item \textbf{片1（高位）}：处理A3=1时，地址1000-1111（译码器输出Y0-Y7对应系统Y8-Y15）
  \item A3=0→片0使能，片1禁止
  \item A3=1→片1使能，片0禁止
\end{itemize}

\textbf{【VHDL】}
\begin{lstlisting}
entity decoder_4to16 is
  port (A : in std_logic_vector(3 downto 0);
        Y : out std_logic_vector(15 downto 0));
end decoder_4to16;

architecture structural of decoder_4to16 is
  signal Y_low, Y_high : std_logic_vector(7 downto 0);
  component decoder_3to8 is
    port (A : in std_logic_vector(2 downto 0);
          E : in std_logic; Y : out std_logic_vector(7 downto 0));
  end component;
begin
  -- 片0：A3=0时使能
  dec_low: decoder_3to8 port map(
    A => A(2 downto 0), E => not A(3), Y => Y_low);

  -- 片1：A3=1时使能
  dec_high: decoder_3to8 port map(
    A => A(2 downto 0), E => A(3), Y => Y_high);

  Y <= Y_high & Y_low;  -- 拼接为16位输出
end structural;
\end{lstlisting}

\textbf{【通用级联公式】}
用$2^k$个$n$-$2^n$译码器构建$(n+k)$-$2^{n+k}$译码器。
高k位地址→片选，低n位地址→各片输入。

\textbf{【常见考题】}用两片74LS138构建4-16译码器。利用使能端做片选。

\end{detailbox}


\section{变式4：用译码器做地址译码}

\begin{detailbox}[完整教学]

\textbf{【应用场景】}
CPU有16位地址总线，但连接了多个外设（RAM、ROM、IO端口等）。
用译码器将地址空间划分为不同区域，产生各外设的片选信号。

\textbf{【例题】}用3-8译码器产生8个外设的片选信号。地址A15-A13做译码输入。

\textbf{【电路】}
\begin{center}
\begin{tabular}{|c|c|l|}
\hline
A15A14A13 & 选中输出 & 外设 \\
\hline
000 & Y0 & RAM1(0x0000-0x1FFF) \\
001 & Y1 & RAM2(0x2000-0x3FFF) \\
010 & Y2 & ROM (0x4000-0x5FFF) \\
011 & Y3 & IO1 (0x6000-0x7FFF) \\
100 & Y4 & IO2 (0x8000-0x9FFF) \\
\hline
\end{tabular}
\end{center}

\textbf{【核心思想】}高位地址→译码→片选。低位地址→片内寻址。

\end{detailbox}


\section{变式5：显示译码器（BCD→7段）}

\begin{detailbox}[完整教学]

\textbf{【问题】}3-8译码器每次只选中1个输出。但7段数码管需要同时点亮多个段（如数字"8"需要7段全亮）。

\textbf{【区别】}
\begin{itemize}
  \item 3-8译码器：$n$输入→$2^n$输出，输出是\textbf{互斥}的（恰好一个为1）
  \item BCD-7段译码器：4输入→7输出，输出\textbf{不互斥}（可以同时多个为1）
\end{itemize}

\textbf{【详见第三章BCD译码器】}BCD-7段译码器本质上是7个独立的组合逻辑函数，不是简单的"最小项展开"。

\begin{examfocus}
\textbf{关键区别}：
\begin{itemize}
  \item 3-8译码器：1-of-N选择（互斥输出）→ 用AND门
  \item BCD-7段：任意多段同时亮（非互斥）→ 用ROM/查找表
\end{itemize}
不搞混这两种"译码器"是考试的基本要求。
\end{examfocus}

\end{detailbox}


\section{译码器变式总结}

\begin{examfocus}[译码器5大变式速查]
\begin{center}
\begin{tabular}{|l|l|l|}
\hline
\textbf{变式} & \textbf{核心考点} & \textbf{常考题型} \\
\hline
74LS138风格 & 低有效+3使能端 & 写VHDL, 判断使能条件 \\
译码器+OR≈任意函数 & 最小项之和 & 给定函数→画电路 \\
级联构建更大译码器 & 高位地址做片选 & 74LS138×2→4-16 \\
地址译码 & 地址空间划分 & 存储器片选设计 \\
vs BCD-7段 & 互斥vs非互斥输出 & 概念辨析 \\
\hline
\end{tabular}
\end{center}
\end{examfocus}
